\documentclass[11pt]{article}
\usepackage[T1]{fontenc}
\usepackage{lmodern,microtype,amsmath,amssymb,amsthm,mathtools,mathrsfs,tikz-cd}
\usepackage[margin=1in]{geometry}
\usepackage{enumitem,booktabs,longtable,xurl,hyperref}
\hypersetup{colorlinks=true,linkcolor=blue,citecolor=blue,urlcolor=blue,
 pdftitle={Iteration domains, finite proofs, and computational cosheaves},pdfauthor={Florian Lengyel}}
\newtheorem{theorem}{Theorem}[section]
\newtheorem{proposition}[theorem]{Proposition}
\newtheorem{lemma}[theorem]{Lemma}
\newtheorem{corollary}[theorem]{Corollary}
\theoremstyle{definition}\newtheorem{definition}[theorem]{Definition}
\theoremstyle{remark}\newtheorem{remark}[theorem]{Remark}
\newcommand{\A}{\mathscr A}
\newcommand{\E}{\mathscr E}
\newcommand{\N}{\mathbb N}
\newcommand{\ran}{\operatorname{ran}}
\newcommand{\Tot}{\operatorname{Tot}}
\newcommand{\Par}{\operatorname{Par}}
\newcommand{\It}{\operatorname{It}}
\newcommand{\Set}{\mathbf{Set}}
\newcommand{\Mono}{\operatorname{Mono}}
\newcommand{\op}{\mathrm{op}}
\title{Iteration domains, finite proofs,\\and computational cosheaves}
\author{Florian Lengyel}
\date{29 October 2026}
\begin{document}
\maketitle
\begin{abstract}
A finite worklist recognizes precisely the codes generated by a constructor
of bounded finite arity in finitely many stages. We assume a split restriction
category with product-compatible ranges, weak choice, a finitely complete
extensive total category, and Heller's stable finite-length semigroup
coverage. Finite execution traces then construct the termination domain.
The resulting union is stable
under parameters, inverse images, and the operations of the inference
polynomial. Weak choice also makes the total category regular, with
computational ranges as regular images. Finite equality constraints in
proof trees give a consistency argument without decidable atom equality.
These results apply to a recursion category of computational cosheaves on
the Sierpi\'nski space. Its universal object has an explicit Turing
morphism and lacks decidable equality. This supplies a nonnumerical
instance of the finite-proof essential-incompleteness construction.
\end{abstract}
\tableofcontents

\section{Iteration structure and computational domains}\label{sec:structure}

The finite-proof construction of \cite{FiniteProofs} requires the union
of the domains of its finite constructor stages. An arbitrary solution
of the fold equation does not identify that union. A terminating
worklist does: decoding a constructor replaces the current code by its
finitely many premises, and an empty worklist records successful
termination. We construct the domain of this computation from finite
traces and compare the traces with constructor trees.

Composition is written $gf=g\circ f$. Let $\A$ be a split restriction
category with restriction zeros and binary restriction products
$\boxtimes$. Write $\bar f:X\to X$ for the restriction of
$f:X\to Y$, and call $e=\bar e$ a \emph{domain}. An arrow is total
if $\bar f=1$. Domains are ordered by $e\le d$ when $de=e$.
Their meet is composition. Splitting a domain means choosing a total
map $i_e:X_e\to X$ and a partial inverse $i_e^\dagger:X\to X_e$
with $i_e^\dagger i_e=1$ and $i_ei_e^\dagger=e$.
Put $\E=\Tot(\A)$; its binary products are denoted by $\times$.
We assume that $\E$ has finite limits and is extensive: finite
coproducts are disjoint and stable, and the canonical functor
$(\E/X)\times(\E/Y)\to\E/(X+Y)$ is an equivalence, including
the empty-coproduct condition. Restriction coproducts in $\A$ agree
with these total coproducts.

Ranges have the universal agreement property. Thus $r=\ran f$ is a
domain on the target with $rf=f$, and
\[
 vf=wf\quad\Longrightarrow\quad vr=wr
\]
for all parallel partial arrows $v,w$. We require the range calculus
and product compatibility
\begin{equation}\label{eq:range-laws}
 \ran(df)=d\ran f,\qquad
 \ran(g\ran f)=\ran(gf),\qquad
 \ran(f\boxtimes g)=\ran f\boxtimes\ran g,
\end{equation}
where $d$ is a domain and each composite is typed. Inverse image is
$f^*d=\overline{df}$. A \emph{partial monomorphism} $f:X\to Y$
satisfies $fu=fv\Rightarrow\bar f u=\bar f v$ for all parallel
$u,v:Z\to X$. Heller's weak choice says that every such $f$ has a
section $s:Y\to X$ with $fs=\bar s$ and $fsf=f$.

A family of domains $(d_i)$ on $X$ has \emph{union} $d$ if
$d_i\le d$ and agreement of parallel partial arrows on every $d_i$
implies agreement on $d$. This is Heller's union, not merely an
order-theoretic supremum. It is \emph{stable} when the same property
holds after product with every identity domain.

Heller's $B^\#$-structure is a covariant functor
$\#:\E\to\E$, on arrows $f\mapsto f^\#$, together with seven
natural transformations \cite[\S2, p.~1255]{Heller}:
\begin{gather*}
 m_X:X^\#\times X^\#\to X^\#,
 \quad j_X:X\to X^\#,
 \quad a_X:X^{\#\#}\to X^\#,\\
 l_X:X^\#\to X+(X\times X^\#),
 \quad r_X:X^\#\to(X^\#\times X)+X,\\
 \operatorname{par}_{X,Y}:X\times Y^\#\to(X\times Y)^\#,
 \quad \operatorname{wd}_{X,Y}:(X+Y)^\#\to D_{X,Y},
\end{gather*}
where
\begin{align*}
 D_{X,Y}={}&X^\#+Y^\#+(X^\#\times Y^\#)^\#
 +(Y^\#\times X^\#)^\#\\
 &+Y^\#\times(X^\#\times Y^\#)^\#
 +X^\#\times(Y^\#\times X^\#)^\#.
\end{align*}
The map $m$ is associative; $j,a$ are the unit and multiplication of
a monad; and $a_X$ is a semigroup homomorphism. The maps $l_X,r_X$
are inverse to $[j_X,m_X(j_X\times1)]$ and
$[m_X(1\times j_X),j_X]$, respectively. The parameter law is
$(\pi_2)^\#\operatorname{par}_{X,Y}=\pi_2$.
To specify the last axiom, put $V=X+Y$,
$u=i_X^\#:X^\#\to V^\#$, $v=i_Y^\#:Y^\#\to V^\#$,
$c=a_V(m_V(u\times v))^\#$, and
$d=a_V(m_V(v\times u))^\#$. Then $\operatorname{wd}_{X,Y}$ is
inverse to $[u,v,c,d,m_V(v\times c),m_V(u\times d)]$.
These equations describe formally free semigroups; they do not
postulate their ordinary free-algebra universal property.

For $n\ge1$, let $j_n:X^n\to X^\#$ concatenate $n$ singletons.
Heller's additional iteration hypothesis is that the domains
$\ran j_n$ have stable union $1_{X^\#}$ for every $X$
\cite[\S7]{Heller}. Together with the preceding weak-choice and
range structure, this is his definition of an iteration category:
a $wr^+$-prodominical category whose total $B^+$-structure extends
to such a $B^\#$-structure.
We use its split form with the stated finite limits and extensivity.
Lengyel's later definition instead assumes countable agreement
unions and an iteration operation for total endomorphisms
\cite[\S2.6, pp.~14--15]{Lengyel}. Its corrected union calculus
states product stability separately \cite[\S11.2]{Repairs}.
Here the unions needed for iteration are constructed from the
finite-length coverage in Heller's definition.
The left and right decompositions give total first and last maps
$X^\#\to X$. Repeated left decomposition shows that each $j_n$ is
monic and that the images for distinct lengths are disjoint. On the
length-$n$ image, first and last are the corresponding projections.

Only the total-arrow action of $\#$, concatenation, first and last,
the length injections, and their stable coverage enter the trace
construction. Write $L(X)=1+X^\#$ for words including the empty
word. The decompositions give $L(X)\cong1+X\times L(X)$ and a
total concatenation $L(X)\times L(X)\to L(X)$, with an empty word
as unit. These follow by extending $m_X$ across the empty summands.

An object $U$ is \emph{isotypical} in Heller's sense if
$U\cong U\times U\cong U+U\cong U^\#$. A \emph{frame} chooses
these three isomorphisms, in the direction from $U$, and their
inverses. The formal objects generated by $U$ under $\times,+,\#$
retain their expression trees; their arrows are maps between their
interpretations in the ambient category. Let $\mathscr B$ be
a $B^\#$-subcategory on these formal objects.
A total catalogue $t:V\times U\to U$ in $\mathscr B$
lists the endomorphisms $f:U\to U$ for which
$t(g\times1_U)=f\pi_2$ for some total $g:U\to V$ in
$\mathscr B$. Uniform generation means that a frame and such a
list generate $\mathscr B$ under the $B^\#$-operations.

For total $u:X\to W$ and $v:W\to W+Y$, the Turing development is
\[
 \operatorname{Tur}(u,v)
 =[0,1_Y]\It([v,i_Y],0_W+1_Y)i_Wu:X\to Y.
\]
Heller's Theorem~9.2 constructs a recursion category
$\operatorname{Tur}(\mathscr B)$ from a uniformly generated
isotypical $\mathscr B$ in a locally connected iteration category.
Here a domain $d$ is connected when $d=d_1\vee d_2$ and
$d_1d_2=0$ imply $d_1=0$ or $d_2=0$. Local connectedness means
that agreement on all connected domains determines arrows.
Lengyel calls this property separation. His existence theorems
permit a faithful iteration functor into a separated category:
it preserves zeros, coproducts, domains, and iteration, but need
not preserve products \cite{Lengyel,Repairs}.
These statements use a fixed frame; no independence of the frame
is needed. Their hypotheses concern the
ambient iteration category. The generated recursion subcategory
requires its own range analysis: Heller's simplicial example already
has a diagonal without a range \cite[\S10, Example~5]{Heller}.

\section{Weak choice, regular images, and stable unions}\label{sec:regular}

The range factorization and its pullback stability appear in
Lengyel's range representation theorem and its corrected proof
\cite[Theorem~42]{Lengyel}, \cite[Appendix~B]{Repairs}.
We give the extra descent argument that turns this factorization
into a regular image factorization under weak choice.
A regular epimorphism is a coequalizer of some pair. A category is
regular if it has finite limits, regular-epimorphism--monomorphism
factorizations, and pullback-stable regular epimorphisms.

\begin{lemma}\label{lem:pullback-range}
Suppose $\A$ is split, has universal ranges satisfying
\eqref{eq:range-laws}, and $\E=\Tot(\A)$ has finite limits.
If $f:X\to Y$ and $g:Z\to Y$ are total and
$p:X\times_Y Z\to Z$ is the pullback of $f$, then
\[
 \ran p=g^*\ran f.
\]
In particular, pullback preserves total arrows with full range.
\end{lemma}
\begin{proof}
The graph $j=\langle g,1_Z\rangle:Z\to Y\times Z$ has range
$G$ and restriction inverse $j^\dagger=\pi_2G$. Indeed,
$\pi_2j=1$ and universal range agreement extends $j\pi_2j=j$
to $j\pi_2G=G$; the restriction of $\pi_2G$ is $G$.
Put $F=f\times1_Z$ and $E=F^*G$. Splitting $E$ gives the
pullback of $f$ and $g$. The restriction identity $FE=GF$ yields
\[
 \ran(FE)=G\ran F=G(\ran f\boxtimes1_Z).
\]
Under the graph isomorphism, the domain on the right is
$j(g^*\ran f)j^\dagger$. Applying $\pi_2$ and taking ranges
therefore gives $\ran p=g^*\ran f$. The use of the splitting
does not change the range, since its inclusion has range $E$.
\end{proof}

\begin{theorem}\label{thm:regular}
Let $\A$ satisfy the hypotheses of Lemma~\ref{lem:pullback-range}
and weak choice. Then every monomorphism in $\E$ is a restriction
monomorphism, $\E$ is regular, and restriction ranges of total maps
are their regular images.
\end{theorem}
\begin{proof}
A total monomorphism $m$ in $\E$ is monic also against partial
arrows. If $mu=mv$, totality gives $\bar u=\bar v$; split this
common domain and use monicity on the resulting total maps.
Weak choice supplies a section $s$ of $m$. From $msm=m$ and
monicity, $sm=1$; hence $ms=\bar s=\ran m$. Thus $m$ is a
restriction monomorphism.

Factor a total $f:X\to Y$ through its split range as
$X\xrightarrow{q}I\xrightarrow{i}Y$. The arrow $q$ is total
and has full range. Such an arrow is epic by universal range
agreement. It remains to prove that it coequalizes its kernel pair.

Suppose $h:X\to Z$ equalizes the kernel pair of $q$. Factor
$\langle q,h\rangle$ through its range as
$X\xrightarrow{d}G\xrightarrow{m}I\times Z$, and write
$u=\pi_1m$, $v=\pi_2m$. The arrows $d$ and $u$ have full range.
For total $a,b:T\to G$ with $ua=ub$, pull $d$ back along $a$
and $b$ and take the fiber product of these two pullbacks over $T$.
The resulting map $z:T'\to T$ has full range by
Lemma~\ref{lem:pullback-range}. There are lifts $x,x':T'\to X$
with $dx=az$ and $dx'=bz$. Since $qx=qx'$, the kernel-pair
hypothesis gives $hx=hx'$, and hence $vaz=vbz$. Epicity of $z$
gives $va=vb$. Together with $ua=ub$ and monicity of $m$, this
gives $a=b$. Thus $u$ is monic.

By the first paragraph, $u$ has a restriction inverse. Its full
range makes that inverse total, so $u$ is an isomorphism. The map
$vu^{-1}:I\to Z$ is the unique map with $(vu^{-1})q=h$.
Thus $q$ is the coequalizer of its kernel pair. Every total arrow
has the required regular image factorization. Conversely, a regular
epimorphism has full range: in its displayed factorization the monic
factor must be invertible, by the coequalizer universal property.
Lemma~\ref{lem:pullback-range} now proves pullback stability.
\end{proof}

Alternatively, suppose that the same range and finite-limit
hypotheses hold and that Heller's division axiom holds: maps with
the same restriction admit a quotient on the range of the divisor
whenever its equality kernel is contained in the other's.
Then the conclusion of Theorem~\ref{thm:regular} follows directly.
For a total monomorphism $m:X\to Y$, divide $1_X$ by $m$;
the quotient is a restriction inverse, because $m$ has the
diagonal as kernel. For a full-range $q:X\to I$ and a map $h$
equalizing its kernel pair, divide $h$ by $q$. The quotient is
total and is the unique descent of $h$. Thus $q$ is again the
coequalizer of its kernel pair, and Lemma~\ref{lem:pullback-range}
proves its stability. This is the factoriality alternative to
weak choice, with the finite limits and ranges retained explicitly.

The same proof gives Heller's factoriality with explicit types.
For $\phi:W\to X$ and $\psi:W\to Y$ with equal restrictions,
assume
\[
 (\phi\boxtimes\phi)^*\ran\Delta_X
 \le(\psi\boxtimes\psi)^*\ran\Delta_Y.
\]
Split their common domain and the range of $\phi$. The inequality
says that the resulting total map to $Y$ equalizes the kernel pair
of the full-range map to the split range of $\phi$. Its descended
map, followed by the partial inverse of the range inclusion, is
$\psi/\phi:X\to Y$, with restriction $\ran\phi$ and
$(\psi/\phi)\phi=\psi$. This derives the factorization rather
than using the ill-typed final composite printed in Heller's
Proposition~6.2.

Finite unions already exist under these hypotheses. For finitely
many domains $d_i$ on $X$, take the range of the total copairing
$[i_{d_i}]:\coprod_iX_{d_i}\to X$. The coproduct and range
universal properties give precisely the agreement-union property.
Distributivity of products over these coproducts and
\eqref{eq:range-laws} give stability. The empty case is $0_X$.

\begin{lemma}\label{lem:unions}
In the split range category of Lemma~\ref{lem:pullback-range},
stable unions are preserved by inverse image along total arrows.
They are also preserved by direct image and finite coproducts.
If $d=\bigcup_i d_i$ is any union, then
$ad=\bigcup_i ad_i$ for every domain $a$ on the same object.
\end{lemma}
\begin{proof}
For the meet assertion, apply agreement on $d$ to $va,wa$.
For inverse image along $f:X\to Y$, put
$j=\langle1_X,f\rangle$, $G=\ran j$, and
$j^\dagger=\pi_1G$. Stability followed by the meet assertion gives
\[
 G(1_X\boxtimes d)=\bigcup_i G(1_X\boxtimes d_i).
\]
The domain-interval isomorphism $c\mapsto j^\dagger c j$
preserves agreement unions: an equality after every
$j^\dagger c_i j$ may be composed with $j^\dagger$ to give an
equality after every $c_i$, extended to their union, and then
composed with $j$. The transported domains are $f^*d$ and
$f^*d_i$. This is the graph proof in
\cite[\S7.1, Proposition on inverse images]{Repairs}.

If $v\ran(fe_i)=w\ran(fe_i)$ for every $i$, then
$vfe_i=wfe_i$ for every $i$. Agreement on $e=\bigcup_i e_i$
and then on the range of $fe$ gives
$v\ran(fe)=w\ran(fe)$. This proves preservation by direct
image. Product compatibility of ranges proves stability of the
image union. For finite coproducts, test the required equality on
each injection and use its union property; distributivity proves
stability.
\end{proof}

We shall also compose covers. If domains $d_i$ have union $1_X$
and each split object $X_{d_i}$ has a domain cover, the images of
the latter covers have union $1_X$: first extend agreement on each
$X_{d_i}$, then on $X$. Products and inverse images of these
composite covers are stable by the lemma.

For increasing families, finite products of their unions are unions
of their finite products with a common index. First take the union
in one coordinate at a time, using stability and the meet assertion;
then bound a finite tuple of indices by its maximum. Thus, for a
finitary inference polynomial, product, inverse-image, and coproduct
preservation are all available under the present range hypotheses.

\section{Finite execution traces}\label{sec:traces}

The next theorem reconstructs Heller's iteration lemma
\cite[Lemma~7.1]{Heller} in the stated split setting.
The formally free semigroup functor acts on total maps. To iterate
a partial map, we apply that functor to its total domain inclusion
and its total value map. This constructs traces without applying
$\#$ to an undefined partial-arrow action.

\begin{theorem}\label{thm:iteration}
Let $\A$ satisfy the structural hypotheses of
Section~\ref{sec:structure}. For every partial endomorphism
$f:X\to X$ and domain $\epsilon$ with $f\epsilon=\epsilon$,
the domains $b_n=\overline{\epsilon f^n}$, $n\ge0$, have a
stable union $b$. There is a unique arrow $\It(f,\epsilon):X\to X$
with restriction $b$ such that
\[
 \It(f,\epsilon)b_n=f^n b_n,
 \qquad \epsilon\It(f,\epsilon)=\It(f,\epsilon).
\]
The construction commutes with restriction products and finite
coproducts of the given data.
\end{theorem}
\begin{proof}
Split $\bar f$ as $i:D\to X$ and put $a=fi:D\to X$; both
are total. Define total maps $A,B:D^\#\to X^\#$ by
\[
 A=m_X\langle j_Xi\,\operatorname{first}_D,a^\#\rangle,
 \qquad
 B=m_X\langle i^\#,j_Xa\,\operatorname{last}_D\rangle.
\]
Let $R\hookrightarrow D^\#$ be their equalizer, further
restricted by $(a\operatorname{last}_D)^*\epsilon$.
Its points, when points are available, are nonempty sequences of
defined steps whose successive endpoints match and whose final
endpoint is fixed by $f$. All objects and maps here are defined by
finite limits and split domains.

Adjoin the zero-step traces by $T=X_\epsilon+R$. The start and
end maps $s,t:T\to X$ are $i_\epsilon$ on the first summand
and $i\operatorname{first}_D$, $a\operatorname{last}_D$ on
the second. Pulling the stable length cover back to $R$ gives
finite-length trace objects $T_n$, with $T_0=X_\epsilon$.
For $n\ge1$, their defining equations on $D^n$ are
\[
 ad_r=id_{r+1}\quad(1\le r<n),
 \qquad ad_n\text{ factors through }X_\epsilon.
\]
The word identities show this by comparing the $n+1$ coordinates
of $Aj_n$ and $Bj_n$. The start map $s_n:T_n\to X$ identifies
$T_n$ with $X_{b_n}$. In one direction, successive application of
$f$ follows these equations. In the other, the total maps
$i^\dagger f^r:X_{b_n}\to D$, $0\le r<n$, give the inverse
tuple. Their totality follows from definedness of $f^n$ there.
Consequently $\ran s_n=b_n$ and $t_n=f^n s_n$.

By Lemma~\ref{lem:unions}, $b=\ran s$ is the stable union of
the $b_n$. To descend $t$ along $s$, take its kernel pair. The
products of the finite-length trace covers and their pullbacks
cover this kernel pair. On its $(n,m)$ component the two start
maps agree. For $k\ge n,m$, the equation $f\epsilon=\epsilon$
gives $t_n=f^k s_n$ and $t_m=f^k s_m$, so the two end maps
agree as well. The cover therefore proves equality on the whole
kernel pair. Theorem~\ref{thm:regular} descends $t$ to the split
range $X_b$ of $s$. Composing with $i_b^\dagger$ gives the
claimed partial map, and the union property gives uniqueness.

For two computations, the $n$th termination domain of their product
is $b_n\boxtimes b'_n$. Independent finite indices are bounded
by their maximum. Lemma~\ref{lem:unions} therefore identifies
the union with $b\boxtimes b'$, and equality on those stages
identifies the maps. The coproduct argument tests each summand.
\end{proof}

The construction uses the finite-length covers, their inverse images,
and their ranges. It requires no externally indexed coproduct of all
finite trace objects.

\section{Constructor stages and a finite worklist}\label{sec:worklist}

Fix $S\in\E$. The partial slice $\A_S$ has total maps
$q:X\to S$ as objects and partial arrows $f:X\to Y$ with
$q'f=q\bar f$ as morphisms. Its total category is $\E/S$.
We call a presentation bounded-finitary if it has the following
form for a fixed natural number $k$.
For $0\le i\le k$, let $A_i$ be a rule-instance object, with
total premise and conclusion maps $s_i:A_i\to S^i$ and
$t_i:A_i\to S$. The covariant inference polynomial is
\[
 P:\E/S\to\E/S,\qquad
 P(q)=\left(\coprod_{i=0}^k A_i\times_{S^i}X^i
       \xrightarrow{[t_i\pi_{A_i}]}S\right).
\]
It acts on a slice map by fixing the rule instance and applying the
map to every premise. Theorem~\ref{thm:regular} makes all total
monomorphisms restriction monomorphisms. Products, pullbacks, and
extensive coproducts therefore extend this action to the domain
spans of partial arrows, giving a restriction functor
$P:\A_S\to\A_S$.

Suppose $q:W\to S$ has a total constructor
$\gamma:Pq\to q$ and a restriction inverse $\delta:q\to Pq$:
\[
 \delta\gamma=1_{Pq},\qquad\gamma\delta=\bar\delta.
\]
As in \cite[Theorem~2.1]{FiniteProofs}, put
\begin{equation}\label{eq:stages}
 e_0=0_W,\qquad e_{n+1}=\gamma P(e_n)\delta.
\end{equation}
These are increasing restriction domains. Let $Q=L(W)$, let
$\eta:W\to Q$ insert a singleton, and let $\epsilon$ be the
empty-word domain. Define a partial endomorphism $F:Q\to Q$
using $Q\cong1+W\times Q$. It fixes the empty word. On a
nonempty word, apply $\delta$ to its head. On the arity-$i$
summand, replace that head by its ordered $i$ premises and append
the unchanged tail. The replacement map
\[
 (A_i\times_{S^i}W^i)\times Q\longrightarrow Q
\]
is total: it discards the rule label and concatenates the finite
premise tuple with the tail. Thus $F$ is a specified partial arrow
and $F\epsilon=\epsilon$.

\begin{theorem}\label{thm:stages}
Under the structural hypotheses of Section~\ref{sec:structure},
the sequence \eqref{eq:stages} has the stable union
\[
 e=\eta^*\overline{\It(F,\epsilon)}.
\]
Moreover, $P(e)$ is the union of $P(e_n)$. Both statements are
stable under adjoining a fixed parameter to every rule, premise,
and conclusion.
\end{theorem}
\begin{proof}
Write $r_m=\eta^*\overline{\epsilon F^m}$. By
Theorem~\ref{thm:iteration} and Lemma~\ref{lem:unions}, $e$ is
the stable union of the $r_m$. We compare finite covers of these
domains with finite covers of the $e_n$.

A constructor shape is a finite ordered rooted tree with arities in
$\{0,\ldots,k\}$. Its height counts vertices along a longest
root-to-leaf path; its size is its number of vertices. Define its
typed label carrier $C_T$ recursively. For a root of arity $i$
and subtrees $T_1,\ldots,T_i$, put
\[
 C_T=A_i\times_{S^i}\prod_{j=1}^i C_{T_j}.
\]
The maps to $S^i$ are the child conclusions. Its code map
$c_T:C_T\to W$ applies $\gamma$ to the rule and the recursively
coded children. At a leaf $C_T=A_0$. Induction gives totality
and monicity of $c_T$. Finite distribution across the arity
summands and \eqref{eq:stages} give
\begin{equation}\label{eq:height-cover}
 e_n=\bigvee_{\operatorname{height}(T)\le n}\ran c_T.
\end{equation}
There are finitely many shapes in this join. For the induction step,
pull the finite child covers back along the premise map and
distribute their finite products; the resulting summands are exactly
the displayed carriers of height at most $n+1$.

Similarly,
\begin{equation}\label{eq:size-cover}
 r_m=\bigvee_{|T|\le m}\ran c_T.
\end{equation}
Here is a finite categorical verification. Unfold the $m$ copies of
$F$ on a singleton input, pulling back the empty/nonempty and arity
coproducts at every step. Extensivity gives a finite cover of the
domain of the resulting composite by these branches. On a branch
with nonempty worklist, the length changes from $l$ to $l-1+i$.
A branch reaching length zero after $j\le m$ nonempty steps is
therefore the prefix traversal of a unique ordered tree with $j$
vertices. The equation $\gamma\delta=\bar\delta$ reconstructs
the head at each step. Reversing the $j$ steps reconstructs the
initial code as $c_T$, with exactly the premise-matching pullbacks
in $C_T$. This gives the upper inclusion in \eqref{eq:size-cover}.
Conversely, $\delta\gamma=1$ makes the worklist traverse the
vertices of $c_T$ in prefix order and become empty after exactly
$|T|$ nonempty steps. It stays empty afterwards, proving the lower
inclusion. These are factorizations of the finite branch maps;
no assumption that points detect arrows is used.

Set $B_0=0$ and $B_{n+1}=1+kB_n$. A tree of height at most $n$
has at most $B_n$ vertices, whereas a tree with at most $m$
vertices has height at most $m$. Equations
\eqref{eq:height-cover}--\eqref{eq:size-cover} give
\[
 e_n\le r_{B_n},\qquad r_m\le e_m.
\]
Thus the two increasing families have the same upper bounds and
the same agreement property. This proves the asserted union of
the $e_n$ and its stability. The final paragraph of
Section~\ref{sec:regular} applies to the finite products, premise
pullbacks, and finite coproduct defining $P$, proving the assertion
about $P(e)$.

For a copied parameter object $K$, use the state object
$K\times Q$ and step $1_K\boxtimes F$. The parameter remains
unchanged when a node is decoded. Its termination domains are
$1_K\boxtimes\overline{\epsilon F^m}$, and its entry map is
$1_K\times\eta$. Product compatibility of iteration gives the
stated parameterized unions.
\end{proof}

\begin{corollary}\label{cor:initial}
Assume the structural hypotheses of Section~\ref{sec:structure}.
Suppose every object admits a total restriction embedding into $U$,
and $\tau:U\boxtimes U\to U$ is parameter-universal: each
$G:Y\boxtimes U\to U$ has a total $c:Y\to U$ with
$\tau(c\boxtimes1_U)=G$. Every bounded-finitary typed rule
presentation then has a parameter-stable initial algebra in $\E/S$.
\end{corollary}
\begin{proof}
For $W=S\times U$, embed the carrier of $Pq$ into $U$ by $k$
and use $\gamma=\langle t,k\rangle$. Restrict its partial
retraction to $\ran\gamma$ to obtain the typed inverse $\delta$,
as in \cite[Proposition~3.2]{FiniteProofs}. Theorem~\ref{thm:stages}
supplies the two unions in the finite-stage initiality theorem.
The typed Turing-fold theorem of \cite[Theorem~4.2]{FiniteProofs}
supplies the partial folds on their exact union. Its evaluator is
filtered by $\ran\Delta_S$ before $P$ is applied, and its strength
identity uses a total program-parameter map. These hypotheses apply
also to copied parameter presentations, proving stability.
\end{proof}

\section{Consistency from finite equality constraints}\label{sec:consistency}

Decidable atom equality was used in \cite[Theorem~6.2]{FiniteProofs}
to turn a finite list of positive axioms into a Boolean valuation.
The finite-stage construction permits a different argument. A fixed
finite proof has only finitely many atom occurrences. Its required
formula equalities impose a finite equivalence relation on those
occurrences. Classical propositional soundness then detects a
positive and a negative axiom in the same equivalence class.

Let \(U\) be the atom object. Use the initial algebra of the covariant
functor
\[
 F_U:\E\to\E,\qquad F_U(X)=U+X+X^2,\qquad
 F_U(f)=1_U+f+(f\times f).
\]
Write \(S\) for its carrier and
\(\theta:U\to S\), \(\nu:S\to S\), and
\({\Rightarrow}:S^2\to S\) for the atom, negation, and implication
constructors. For a domain \(d\) on \(S\), the inverse image
\(\nu^*d=\overline{d\nu}\) is monotone in \(d\); it is a
domain of codes, not Boolean complementation of \(d\).
Initiality makes the structure map \(U+S+S^2\to S\)
an isomorphism: the algebra map to its image under \(F_U\), and
initiality of the original algebra, give the two inverse equations.
Thus these constructors are monic and their images are disjoint.

Let \(a,b\) be disjoint domains on \(U\), with split inclusions
\(m_A:A\to U\) and \(m_B:B\to U\). The positive and negative atomic
axioms have conclusion maps \(\theta m_A\) and \(\nu\theta m_B\).
The logical schemes, with \(\alpha,\beta,\gamma\) ranging over \(S\),
are
\begin{align*}
 H_1(\alpha,\beta)&=\alpha\Rightarrow(\beta\Rightarrow\alpha),\\
 H_2(\alpha,\beta,\gamma)&=
  (\alpha\Rightarrow(\beta\Rightarrow\gamma))
  \Rightarrow((\alpha\Rightarrow\beta)
                 \Rightarrow(\alpha\Rightarrow\gamma)),\\
 H_3(\alpha,\beta)&=
  (\nu\beta\Rightarrow\nu\alpha)
                 \Rightarrow(\alpha\Rightarrow\beta).
\end{align*}
There is one binary rule: from \(\alpha\) and
\(\alpha\Rightarrow\beta\), conclude \(\beta\).
These data define the slice polynomial with nullary carrier
\(A+B+S^2+S^3+S^2\) and binary carrier \(S^2\). Its object and
morphism actions are those in Section~\ref{sec:worklist}.
Let \(c:D\to S\) be the initial proof algebra constructed there,
and let \(t=\ran c\).

\begin{lemma}\label{lem:shapes}
The syntax and proof objects just constructed have stable covers by
finite constructor shapes. A proof shape can be further covered by
choosing finite syntax shapes for all its formula parameters.
For any fixed pair of such choices, equality of two displayed
formulas either has zero domain or is a finite conjunction of
equalities between atom labels.
\end{lemma}
\begin{proof}
Equation~\eqref{eq:height-cover} and the stage union give the syntax
cover. Applying the same equation to the rule polynomial gives
the proof cover. First refine each nullary label by the finite
coproduct \(A+B+S^2+S^3+S^2\). Extensivity gives a finite cover
on which each leaf has a fixed signed-axiom or logical-scheme tag.
A fixed proof tree then has finitely many rule labels,
each using finitely many formula parameters. Products of their
syntax covers, followed by the pullbacks imposing premise matching,
give the refinement, by Lemma~\ref{lem:unions}. These choices need
not have a common height bound across the whole proof object.

For the last assertion, compare two fixed formula trees at the root.
Different constructor tags have disjoint images. Equal negation
tags reduce to their arguments; equal implication tags reduce to
both pairs of arguments; equal atom tags reduce, by monicity of
\(\theta\), to equality in \(U\). Finite induction proves the
assertion as an equality of subobjects. All equality subobjects
are represented by restriction domains by Theorem~\ref{thm:regular}.
\end{proof}

\begin{lemma}[Finite equality obstruction]\label{lem:finite-obstruction}
Fix two finite Hilbert proof trees whose conclusions are required
to be a formula and its negation. Fix formula shapes for their
parameters, and impose all equalities required by their rules and
conclusions. If these constraints do not already have zero domain,
their finite equality relation identifies the label of a positive
atomic axiom with the label of a negative atomic axiom.
\end{lemma}
\begin{proof}
By Lemma~\ref{lem:shapes}, all remaining matching constraints are
equalities between a finite set of atom occurrences. Let \(\sim\)
be the equivalence relation generated by these required equalities.
Assign a propositional variable to each equivalence class. This
turns each tree into a finite propositional derivation from its
signed atomic leaves. Repeated arguments of a scheme and matched
premises have the same variables, because their occurrence
equalities were included.

If no class contains both a positive and a negative atomic leaf,
assign true to the classes containing positive leaves and false to
all remaining classes. Every signed leaf is true. The three
displayed schemes are true by their Boolean truth tables, and
modus ponens preserves truth. Both conclusions would therefore
be true, although they are Boolean complements. This contradiction
proves that a class has both kinds of leaf. Its equality chain is
a finite chain of the imposed atom equalities.
\end{proof}

\begin{theorem}\label{thm:consistency}
Assume the hypotheses of Corollary~\ref{cor:initial}. For disjoint
domains \(a,b\) on \(U\), the finite-stage syntax and proof algebras
defined above satisfy
\[
 t\wedge\nu^*t=0_S,\qquad t=\ran(c:D\to S).
\]
Decidable equality on \(U\) is not required.
\end{theorem}
\begin{proof}
Form the object \(Q\) of proof/refutation pairs by the pullback
\[
\begin{tikzcd}
 Q\arrow[r,"p_-"]\arrow[d,"p_+"']&D\arrow[d,"c"]\\
 D\arrow[r,"\nu c"']&S.
\end{tikzcd}
\]
Its two conclusions satisfy \(cp_-=\nu cp_+\).
The finite proof covers, their syntax refinements, their products,
and the pullback in this diagram give a stable cover of \(Q\).
Consider one member \(Z\to Q\). A constructor mismatch makes
\(Z\cong0\). Otherwise Lemma~\ref{lem:finite-obstruction} gives
a positive leaf map \(x:Z\to A\) and a negative leaf map
\(y:Z\to B\) with \(m_Ax=m_By\). Hence \(Z\) maps to
\(A\times_U B\cong0\). The initial object is strict by
extensivity, so again \(Z\cong0\).

All domains in the cover are consequently zero. Their agreement
union is \(1_Q\), so \(1_Q=0_Q\) and \(Q\cong0\).
Pullback stability of regular images, including their agreement
with restriction ranges, now gives \(t\wedge\nu^*t=0_S\).
\end{proof}

The proof uses only the equalities required by a fixed finite
derivation. Additional equalities among its actual atom labels do
not affect the argument: the derived equality between an \(A\)-leaf
and a \(B\)-leaf still holds. No enumeration of disequality cases
or Boolean complement of an equality domain is needed.

For comparison, the finite partition criterion can be stated without
a categorical completion. Given a fixed formula with \(n\) atom
occurrences, a partition \(\pi\) of those occurrences identifies
variables within its blocks. Select positive and negative literal
subsets \(I,J\). Retain the triple when
\[
 \left(\bigwedge_{i\in I}x_{[i]_\pi}\ \wedge\
       \bigwedge_{j\in J}\neg x_{[j]_\pi}\right)
       \ \Longrightarrow\ \varphi_\pi
\]
is a Boolean tautology. Its domain is a finite meet of within-block
equality domains and inverse images of \(a,b\); the finite join
over retained triples exists. For ordinary sets this characterizes
semantic consequence from the signed atomic axioms. One direction
is substitution into a tautology. For the other, partition by the
actual equalities and assign the remaining finitely many variables
freely. This criterion and Lemma~\ref{lem:finite-obstruction} use
positive equality information. The categorical proof above obtains
its selected theorem image directly as \(\ran c\).

\begin{corollary}\label{cor:essential}
Under the hypotheses of Corollary~\ref{cor:initial}, suppose \(U\)
is nonzero and the Turing morphism has the binary case distinction
of \cite[Proposition~1.1]{FiniteProofs}. The resulting inseparable
pair \(a,b\) gives a consistent computational presentation such that
every consistent computational extension is incomplete.
\end{corollary}
\begin{proof}
Theorem~\ref{thm:consistency} gives consistency. Its theorem domain
is the least rule-closed domain: factor \(c\) through its regular
image and use preservation of regular epimorphisms by the finite
products, pullbacks, and coproducts defining \(P\). The proof
algebra structure descends to that image, and initiality maps into
every other rule-closed domain.

The one-node proof maps give
\(a\le\theta^*t\) and \(b\le\theta^*\nu^*t\).
For any computational domain \(t'\ge t\), consistency means
\(t'\wedge\nu^*t'=0\), and completeness would mean
\(t'\vee\nu^*t'=1\). Pulling these equations back along the total
map \(\theta\) would separate \(a,b\) by complementary domains.
This contradicts their inseparability. A computational theory is
a computational domain closed under the displayed schemes and rule.
\end{proof}

\section{A computational cosheaf recursion category}\label{sec:cosheaves}

Let \(X\) be a topological space, and let \(\mathcal O(X)\) be the
category of open sets and inclusions. A set-valued cosheaf is a
covariant functor \(C:\mathcal O(X)\to\Set\) for which the diagram
\[
 \coprod_{i,j}C(V_i\cap V_j)
 \rightrightarrows\coprod_i C(V_i)\longrightarrow C(V)
\]
is a coequalizer for every open cover \(V=\bigcup_iV_i\), with
\(C(\varnothing)=\varnothing\). Its arrows are natural
transformations, composed componentwise; the two parallel arrows
come from the inclusions of intersections.
This fixes the variance used here \cite[\S\S2.1--2.2]{Curry}.

Take the Sierpi\'nski space \(X=\{0,1\}\) with opens
\(\varnothing,V=\{1\},X\). Each cover of a nonempty open contains
that open itself. Thus a cosheaf is exactly a map
\[
 r_C:C(V)\longrightarrow C(X).
\]
An arrow between cosheaves is a commuting square. Evaluation on
the two nonempty opens is an equivalence
\[
 \operatorname{CoSh}(X,\Set)\ \simeq\ \Set^{[1]},
\]
where \([1]\) has two objects and one nonidentity arrow \(0\to1\).
The inverse functor puts the empty set at \(\varnothing\) and
the given arrow at \(V\subset X\). Identities, composition, and
both inverse identifications follow componentwise. In particular,
a noninvertible arrow is allowed.

Write c.e.\ for computably enumerable.
Let \(\mathscr C\) be the category of c.e. subsets of finite powers
of \(\N\), with recursive functions total on their sources as
arrows. Write
\[
 \E_\to=\mathscr C^{[1]},\qquad
 \A_\to=\Par(\E_\to,\Mono).
\]
We call its total objects computational cosheaves. Their arrows
are the computational commuting squares, rather than all maps of
their underlying sets. A partial arrow is a commuting square
defined on a c.e. subdiagram. Explicitly, for
\(r_X:X_0\to X_1\) and \(r_Y:Y_0\to Y_1\), it consists of c.e.
subsets \(D_i\subseteq X_i\) with \(r_X(D_0)\subseteq D_1\) and
partial recursive value maps \(f_i:D_i\to Y_i\), total on \(D_i\),
such that \(r_Yf_0=f_1r_X|_{D_0}\). Equality means equality of
the two domain subsets and their value maps. Composition takes
componentwise inverse images of domains and composes the values.
Products and coproducts also use the two components.

\begin{proposition}\label{prop:cosheaf-structure}
The category \(\A_\to\) satisfies all structural hypotheses of
Section~\ref{sec:structure}. Its total category is regular and
extensive, and its ranges are the componentwise c.e. images.
\end{proposition}
\begin{proof}
Finite limits of \(\mathscr C\) are given by tuples and recursive
equality tests. Images are c.e. A total surjection has a recursive
section on its target: enumerate its c.e. source and search for
a preimage. Consequently its image factorization is regular and
regular epimorphisms are pullback-stable. Tagged sums are disjoint
and stable. These constructions give the indicated properties of
\(\mathscr C\).

Finite limits, finite coproducts, and images in the arrow category
are computed componentwise. The image transition exists by the
commuting-square equation, and is total recursive on the image.
The objectwise coequalizer property of a regular epimorphism gives
a natural descended map: its naturality follows after
precomposition with the epic source component. Thus the componentwise
image factorization is regular, with pullback-stable regular
epimorphisms. Componentwise coproduct decompositions prove
extensivity.

All monomorphisms in \(\E_\to\) are componentwise injections.
For component \(0\), test against maps from \((1\to1)\), and
for component \(1\), against maps from \((0\to1)\); their arbitrary
chosen values are computable constants. Conversely componentwise
injectivity implies monicity. Such a mono has a partial recursive
inverse on its c.e. componentwise image, so it represents a selected
subdiagram inclusion up to isomorphism.

The partial-map construction is split and has restriction products,
coproducts, and zeros. Regular images supply universal ranges:
equality of two partial composites includes equality of their
pulled-back domains; surjectivity of the image map reflects this
equality and equality of values. Images commute with products.
For a partial monomorphism represented by
\(X\xleftarrow{m}D\xrightarrow{h}Y\), the right leg \(h\) is
monic, by testing total arrows into \(D\). Reversing the span
gives its section. This proves weak choice.

Finally, nonempty words in each c.e. component form a c.e. carrier,
and the transition applies \(r_X\) to every letter. All seven
\(B^\#\)-maps are the recursive word operations, natural under
the transition. The length injections cover every word in each
component. Pointwise agreement, with any fixed parameter retained,
proves their stable union property. This does not assert existence
of a c.e. union for a nonuniform external sequence of c.e. domains.
\end{proof}

The universal object has a noninjective transition:
\[
 U=\bigl(\N^2\xrightarrow{\pi_2}\N\bigr).
\]
For any \(r_Y:Y_0\to Y_1\), choose computable injections
\(j_i:Y_i\to\N\), with partial recursive inverses. The square
\[
\begin{tikzcd}
 Y_0\arrow[r,"{\langle j_0,j_1r_Y\rangle}"]\arrow[d,"r_Y"']&
 \N^2\arrow[d,"\pi_2"]\\
 Y_1\arrow[r,"j_1"']&\N
\end{tikzcd}
\]
commutes and is monic, so it is a total restriction embedding
into \(U\).

\begin{proposition}\label{prop:cosheaf-turing}
The object \(U\) has a parameter-universal Turing morphism
\(\tau:U\boxtimes U\to U\) in \(\A_\to\).
\end{proposition}
\begin{proof}
Fix a universal partial recursive function \((p,x)\mapsto\varphi_p(x)\)
with a total recursive specialization function, and a recursive
pairing bijection \(\N^2\cong\N\). At component \(1\), set
\[
 \tau_1(p_1,x_1)=\varphi_{p_1}(x_1).
\]
At component \(0\), the program and input are pairs. Define
\[
 \tau_0\bigl((p_0,p_1),(x_0,x_1)\bigr)=(y_0,y_1)
\]
exactly when the following two computations both halt and have
the indicated values:
\[
 \varphi_{p_0}(\langle x_0,x_1\rangle)=\langle y_0,y_1\rangle,
 \qquad \varphi_{p_1}(x_1)=y_1.
\]
Dovetailing and recursive equality give this partial recursive map.
Its domain projects into the domain of \(\tau_1\), and its output
projects to \(\tau_1(p_1,x_1)\). Hence it is a partial arrow of
\(\A_\to\), with its stated domain subdiagram.

Now let \(F:Y\boxtimes U\to U\) be an arbitrary partial arrow.
Choose programs for its two components, including the c.e. domain
tests. Specialization supplies total recursive maps \(s_0,s_1\)
with
\[
 \varphi_{s_0(y_0)}(\langle x_0,x_1\rangle)
     =\langle F_0(y_0,(x_0,x_1))\rangle,\qquad
 \varphi_{s_1(y_1)}(x_1)=F_1(y_1,x_1),
\]
as equalities of partial computations. The brackets on the first
output mean the numerical code of its output pair.
Define
\[
 c_1(y_1)=s_1(y_1),\qquad
 c_0(y_0)=(s_0(y_0),s_1(r_Yy_0)).
\]
These maps are total on \(Y_i\), and
\(\pi_2c_0=c_1r_Y\), so \(c:Y\to U\) is total.
At component \(1\), universality gives \(\tau_1(c_1,-)=F_1\).
Where \(F_0\) is defined, the domain and naturality conditions
for \(F\) make the extra component-\(1\) test in \(\tau_0\)
succeed. Conversely, that evaluator cannot halt unless the
specialized \(F_0\) program halts. Thus
\(\tau(c\boxtimes1_U)=F\), including equality of domains.
\end{proof}

There are computable frames at \(U\). For \(U\cong U^2\), use a
pairing bijection on the base coordinate and an independent one
on the fiber coordinate. For \(U\cong U+U\), the parity of the
base coordinate chooses the summand, and the fiber coordinate is
unchanged. For \(U\cong U^\#\), decode the base as a nonempty
word of length \(n\), then use a recursive bijection
\(\N\cong\N^n\) to supply the \(n\) fiber coordinates. The
inverses recode the base and fiber data, and all three squares
commute with their transitions.

Let \(\mathscr R_\to\) be the full partial subcategory of
\(\A_\to\) on the formal objects generated by \(U\) under
\(\times,+,\#\). Its frame makes it isotypical, and
Proposition~\ref{prop:cosheaf-turing} makes it a recursion category.
Every object of \(\A_\to\) embeds into \(U\), so
\(\A_\to\) is equivalent to the restriction-idempotent splitting
of \(\mathscr R_\to\).
The range domains are computed within this category by
Proposition~\ref{prop:cosheaf-structure}. This is an explicit
computational instance of the functor-category setting in
Heller's Example~4, expressed covariantly as cosheaves.

\begin{proposition}\label{prop:no-decidable-equality}
The universal object \(U\) has no decidable equality in \(\E_\to\).
The category \(\E_\to\) is not equivalent to the numerical total
category \(\mathscr C\).
\end{proposition}
\begin{proof}
The Boolean object \(2=1+1\) is the identity arrow \(2\to2\).
A putative equality map \(\epsilon:U\times U\to2\) must satisfy
\[
 \epsilon_0((x_0,x_1),(y_0,y_1))=\epsilon_1(x_1,y_1).
\]
Its component \(1\) is true when \(x_1=y_1\). For \(x_0\ne y_0\)
and \(x_1=y_1\), the displayed equation makes component \(0\)
true as well, contradicting its required diagonal inverse image.
More generally, decidable equality in this arrow category forces
the transition map to be injective.

The terminal object \((1\to1)\) has exactly three subobjects up
to isomorphism:
\[
 (0\to0),\qquad(0\to1),\qquad(1\to1).
\]
The numerical terminal object has two. An equivalence preserves
these subobject orders, proving the second assertion.
\end{proof}

\begin{theorem}\label{thm:cosheaf-essential}
The recursion category \(\mathscr R_\to\), in its restriction
splitting, admits the finite-stage construction of a consistent
computational propositional theory for which every consistent
computational extension is incomplete.
\end{theorem}
\begin{proof}
The structural, dominance, and Turing hypotheses have been
verified above. For binary case distinction choose the two global
points of \(U\) with components \((0,i)\) and \(i\), for
\(i=0,1\), and let \(q_i\) be their range domains.
Given complementary domains \(c,c'\) on \(U\), their component
subsets are complementary c.e. sets and hence decidable by parallel
search. The subdiagram conditions imply that the component-\(0\)
partition is the inverse image of the component-\(1\) partition
along \(\pi_2\): membership in the opposite part would project
into that opposite part.

Define a total square \(h:U\to U\) whose component \(1\) returns
\(0\) on \(c'_1\) and \(1\) on \(c_1\), and whose component \(0\)
returns \((0,h_1\pi_2)\). Then \(h^*q_0=c'\) and \(h^*q_1=c\).
Also \(0_U\ne1_U\). The Turing diagonal therefore supplies the
inseparable pair used in Corollary~\ref{cor:essential}.
That corollary applies despite Proposition~\ref{prop:no-decidable-equality}.
The theory, its proof object, and its selected theorem range all
belong to the computational cosheaf category.
\end{proof}

\section{Cosheaf products and further constructions}\label{sec:limits}

The Sierpi\'nski example has no nontrivial covering condition on a
nonempty open. Its computational content lies in the noninvertible
extension map, the compatible partial domains, and the filtered
Turing evaluator. On spaces with nontrivial covers, even the
polynomial object construction requires a separate check.

For example, on a discrete two-point space \(X=\{x,y\}\), the
cosheaf \(C(V)=V\) has singleton values on singleton opens.
The pointwise product assignment has four elements at \(X\),
whereas the coproduct of its values on \(\{x\}\) and \(\{y\}\)
has two. Thus pointwise products of set-valued cosheaves need
not be cosheaves. The actual categorical product on this discrete
space pairs elements over the same point.

Likewise, the assignment taking an open \(V\) to formulas whose
atoms lie in \(C(V)\) fails cosheaf descent: an implication with
one atom from each singleton exists over \(X\) and cannot come
from either member of the disjoint cover. Internal polynomial
algebras must therefore use the actual products of the cosheaf
category. Their relationship to an existing computational syntax
object must be proved.

Lengyel's corrected coalgebra construction gives another source
of ambient regular and extensive categories, with product-stable
ranges and unions \cite[\S\S8--10]{Repairs}. Its explicit
two-phase example is equivalent to the numerical recursion
category; it does not supply the nonnumerical distinction proved
in Proposition~\ref{prop:no-decidable-equality}. For a further
coalgebra or cosheaf example, the remaining task is to exhibit a
computational subcategory closed under the range and finite-trace
constructions, with universal-object dominance and a Turing
morphism. Ambient ordinary images alone do not establish that
closure.

The next cosheaf case is the three-point space with opens
\(\varnothing,\{c\},\{a,c\},\{b,c\},\{a,b,c\}\).
Its two proper maximal opens cover the whole space. A cosheaf
must realize the pushout of the two extension maps from
\(C(\{c\})\). A computational construction must identify which
such pushouts admit computational presentations, prove the
corresponding product and range laws, and construct a compatible
Turing evaluator. The two-open refinement has no analogue of the
automatic Sierpi\'nski descent argument.

\begin{thebibliography}{9}
\bibitem{FiniteProofs} F.~Lengyel,
\emph{Finite proof algebras in range restriction categories},
29 October 2026, repository version
\texttt{dbd8b1ad803c1436ecf6abdc1c4ef1548f7ab4c1}.
\url{https://github.com/flengyel/CategoricalEssentialIncompleteness/blob/dbd8b1ad803c1436ecf6abdc1c4ef1548f7ab4c1/manuscript/general_recursion_category_construction.tex}.
\bibitem{Heller} A.~Heller,
An existence theorem for recursion categories,
\emph{Journal of Symbolic Logic} \textbf{55}(3) (1990), 1252--1268.
\href{https://doi.org/10.2307/2274486}{doi:10.2307/2274486}.
\bibitem{Lengyel} F.~Lengyel,
More existence theorems for recursion categories,
\emph{Annals of Pure and Applied Logic} \textbf{125} (2004), 1--41.
\href{https://doi.org/10.1016/S0168-0072(03)00060-5}{doi:10.1016/S0168-0072(03)00060-5}.
\bibitem{Repairs} F.~Lengyel,
\emph{Corrections to More existence theorems for recursion categories},
version 10, 20 October 2026, repository version \texttt{5555e4f}.
\url{https://github.com/flengyel/errata/blob/5555e4fa1da7b8e3a0836dd1fbf1ca385cacd03b/lengyel2004/lengyel2004_repairs_v10.tex}.
\bibitem{Curry} J.~M.~Curry,
\emph{Sheaves, Cosheaves and Applications},
Ph.D. thesis, University of Pennsylvania, 2014.
\href{https://arxiv.org/abs/1303.3255}{arXiv:1303.3255v2},
17 December 2014.
\end{thebibliography}
\end{document}


